\section{Korean external benchmarks}
\label{sec:kbench}

The product evaluation of Section~\ref{sec:data} measures the six tasks Bobcat is built for.
To check that task training did not damage general Korean understanding, we evaluate \expone{}
and the same-input zero-shot base model on public benchmarks that were never used for
training, selection or calibration, through the same typed interface. The suite is KoBEST
\citep{kim2022kobest} (BoolQ, COPA, WiC, HellaSwag, SentiNeg), KMMLU \citep{son2024kmmlu},
the multiple-choice tasks of HAE-RAE Bench~1.1 \citep{son2023haerae} and CLIcK
\citep{kim2024click}, with MMLU \citep{hendrycks2021mmlu} and HellaSwag
\citep{zellers2019hellaswag} as English controls. On all four Korean benchmarks \expone{} is
more accurate than the zero-shot model, by 4.6 to 6.9 points, and every 95\% interval lies
above zero. On the English controls the difference is +3.7 points on MMLU and $-0.1$ points
on HellaSwag (no difference). The numbers in this section are from
\texttt{reports/2026-09-25-korean-external-benchmarks.json}.

\subsection{Method}

\paragraph{Data.}
Each benchmark was downloaded from the Hugging Face Hub at a pinned dataset revision
(Table~\ref{tab:kbench_sources}), and every file was checked against the Hub's listing (LFS
SHA-256, or the git blob SHA-1 for files not stored in LFS) before use. We use the public
test splits where they exist: the five KoBEST test sets, the 45 KMMLU test subjects and the
MMLU test set; for HAE-RAE Bench~1.1 all eleven multiple-choice tasks (the two generative
tasks, lyrics and proverb denoising, are excluded); and CLIcK in full. For HellaSwag we
draw 2,000 items from the validation split with a fixed seed. Items whose text key matches a
text key of the \expone{} training mixture (Section~\ref{sec:mixture}, 50,938 rows) were
removed before evaluation: 3 from KoBEST and 1 from MMLU. One HAE-RAE item exceeds the
8,192-token input limit of this run; it counts as a failure, and therefore as wrong, for both
models. The public test splits are development data for Bobcat: nothing about the model
was chosen by looking at them, but they are not a sealed split either.

\begin{table}[t]
\centering
\small
\caption{Benchmark sources. Revisions are the first eight hexadecimal digits of the pinned
Hub dataset revision. Items are counted after removing the four items whose text overlaps
the training mixture.}
\label{tab:kbench_sources}
\begin{adjustbox}{max width=\linewidth}
\begin{tabular}{@{}lllrl@{}}
\toprule
Benchmark & Hub dataset & Revision & Items & Composition \\
\midrule
KoBEST & \texttt{skt/kobest\_v1} & \texttt{a5ea15e3} & 4,558 & test, 5 tasks \\
KMMLU & \texttt{HAERAE-HUB/KMMLU} & \texttt{d61b3f19} & 35,030 & test, 45 subjects \\
HAE-RAE Bench 1.1 & \texttt{HAERAE-HUB/HAE\_RAE\_BENCH\_1.1} & \texttt{b480e810} & 3,606 & 11 multiple-choice tasks \\
CLIcK & \texttt{EunsuKim/CLIcK} & \texttt{d6162785} & 1,995 & 11 categories \\
\midrule
MMLU (English) & \texttt{cais/mmlu} & \texttt{c30699e8} & 14,041 & test, 57 subjects \\
HellaSwag (English) & \texttt{Rowan/hellaswag} & \texttt{218ec52e} & 2,000 & seeded sample of validation \\
\bottomrule
\end{tabular}
\end{adjustbox}
\end{table}

\paragraph{Typed requests.}
Each item becomes one request with one question and no examples. For KMMLU, HAE-RAE, CLIcK,
MMLU and HellaSwag the question is a Choice whose candidates are the benchmark's own options,
named \texttt{A} to \texttt{E}, with the option text as the candidate description; the state
holds the item text (the question with its passage, if any, or the HellaSwag context), and
the instruction is one
short sentence in Korean, or in English for the controls (for KMMLU, ``choose the correct
answer to the following question''). KoBEST keeps its task structure: BoolQ and WiC are Noul
questions (the probability that the answer to the question is yes, and that the word has the
same meaning in both sentences), and COPA, HellaSwag and SentiNeg are Choices over the two
alternatives, the four endings and the two sentiment labels.

\paragraph{Readout and arms.}
The model generates nothing. It reads the logits of the offered candidate identifiers at the
first answer position (Section~\ref{sec:readout}), as in the product path, and the answer is
the argmax (for a Noul, \emph{true} if $P(\text{true})>0.5$). Both arms receive the same
piecewise inputs (Section~\ref{sec:piecewise}) through the evaluation path of the sealed
final (Section~\ref{sec:final}): \expone{} with its LoRA adapter unmerged, and the zero-shot
arm with no adapter. Probabilities use each arm's product calibration temperature, fitted
beforehand on the product calibration split (Section~\ref{sec:calibration}; $T=1.149$ for
\expone{} and $T=1.26$ for the zero-shot model). Nothing was fitted, selected or tuned on these
benchmarks.

\paragraph{Metrics and intervals.}
Accuracy counts failures as wrong. Chance is the mean of $1/K$ over a benchmark's items, $K$
being the number of candidates. ECE uses ten equal-width bins of the top candidate
probability over the scored items, both raw ($T=1$) and at the arm's temperature. The
difference \expone{} minus zero-shot is reported with a paired 95\% percentile interval from
1,000 cluster-bootstrap resamples. A cluster is a group of items that share a passage or an
article (KoBEST, and the HAE-RAE and CLIcK items that come with a passage) or, for HellaSwag,
a source video; every other item is its own cluster.

\paragraph{Compute.}
The run used one node of eight H100 GPUs for about 1.22\,h. The first attempt was cut off by
an automatic operating-system reboot and was restarted.

\subsection{Results}

\begin{table}[t]
\centering
\small
\caption{Public benchmark accuracy (\%) through the typed first-position readout. $\Delta$ is
\expone{} minus zero-shot in points, with a paired 95\% cluster-bootstrap interval. ECE is
given raw and at the arm's product temperature $T$ (raw\,/\,$T$). Failures count as wrong
(one HAE-RAE item, for both arms). MMLU and HellaSwag are the English controls.}
\label{tab:kbench}
\setlength{\tabcolsep}{4pt}
\begin{adjustbox}{max width=\linewidth}
\begin{tabular}{@{}llrrrrlll@{}}
\toprule
 & & & & \multicolumn{2}{c}{Accuracy} & & \multicolumn{2}{c}{ECE (raw\,/\,$T$)} \\
\cmidrule(lr){5-6}\cmidrule(l){8-9}
Benchmark & Lang. & Items & Chance & Zero-shot & \expone{} & $\Delta$ [95\%] & \expone{} & Zero-shot \\
\midrule
KoBEST & ko & 4,558 & 47.3 & 88.3 & \textbf{95.2} & \textbf{+6.9} [+6.0, +7.7] & 0.009\,/\,0.021 & 0.047\,/\,0.077 \\
KMMLU & ko & 35,030 & 25.0 & 58.1 & \textbf{64.1} & \textbf{+6.0} [+5.5, +6.4] & 0.076\,/\,0.049 & 0.068\,/\,0.028 \\
HAE-RAE Bench 1.1 & ko & 3,606 & 20.0 & 66.8 & \textbf{72.3} & \textbf{+5.5} [+4.1, +6.9] & 0.058\,/\,0.030 & 0.021\,/\,0.035 \\
CLIcK & ko & 1,995 & 24.4 & 71.5 & \textbf{76.1} & \textbf{+4.6} [+2.9, +6.1] & 0.045\,/\,0.024 & 0.027\,/\,0.023 \\
\midrule
MMLU & en & 14,041 & 25.0 & 80.6 & 84.3 & +3.7 [+3.2, +4.2] & 0.029\,/\,0.008 & 0.056\,/\,0.023 \\
HellaSwag & en & 2,000 & 25.0 & 93.3 & 93.2 & $-0.1$ [$-1.0$, +0.9] & 0.029\,/\,0.050 & 0.013\,/\,0.036 \\
\bottomrule
\end{tabular}
\end{adjustbox}
\end{table}

\begin{figure}[t]
\centering
\figfile[width=0.85\linewidth]{figures/korean.pdf}
\caption{Accuracy of \expone{} and the same-input zero-shot model on four Korean benchmarks
and two English controls, with the chance level (dotted) and the paired difference in points
with its 95\% cluster-bootstrap interval above each pair.}
\label{fig:korean}
\end{figure}

Table~\ref{tab:kbench} and Figure~\ref{fig:korean} give the results. \expone{} is more
accurate on every Korean benchmark, by +6.9 points on KoBEST, +6.0 on KMMLU, +5.5 on HAE-RAE
and +4.6 on CLIcK, and each interval excludes zero. On the English controls the difference
is +3.7 [+3.2, +4.2] on MMLU and $-0.1$ [$-1.0$, +0.9] on HellaSwag. Raw ECE of \expone{} is
lower than that of the zero-shot model on KoBEST and MMLU and higher on KMMLU, HAE-RAE, CLIcK
and HellaSwag.

Table~\ref{tab:kbench_sub} breaks KoBEST and HAE-RAE down by task. The largest KoBEST gain is
on WiC (90.9\% against 76.6\%), followed by BoolQ (+6.5) and HellaSwag (+5.8); COPA and
SentiNeg, already near 97\% before training, move by about one point. Within HAE-RAE, CSAT
geography (+16.7), CSAT law (+10.6), reading comprehension (+9.6) and general knowledge (+8.0)
improve most, while date understanding does not improve (50.3\% against 51.2\%).

\begin{table}[t]
\centering
\small
\caption{KoBEST and HAE-RAE Bench~1.1 accuracy (\%) by task. Task names for HAE-RAE are the
dataset's own identifiers. $\Delta$ is in points and is computed before rounding, so it can
differ by 0.1 from the difference of the rounded accuracies. No per-task intervals were
computed.}
\label{tab:kbench_sub}
\begin{tabular}{@{}llrrrr@{}}
\toprule
Benchmark & Task & Items & Zero-shot & \expone{} & $\Delta$ \\
\midrule
KoBEST & BoolQ (Noul) & 1,404 & 91.3 & 97.8 & +6.5 \\
 & COPA (Choice) & 1,000 & 97.8 & 98.7 & +0.9 \\
 & WiC (Noul) & 1,257 & 76.6 & \textbf{90.9} & \textbf{+14.3} \\
 & HellaSwag (Choice) & 500 & 84.0 & 89.8 & +5.8 \\
 & SentiNeg (Choice) & 397 & 96.2 & 97.2 & +1.0 \\
\midrule
HAE-RAE & \texttt{reading\_comprehension} & 936 & 74.3 & 83.9 & +9.6 \\
 & \texttt{date\_understanding} & 475 & 51.2 & 50.3 & $-0.8$ \\
 & \texttt{correct\_definition\_matching} & 439 & 85.2 & 86.1 & +0.9 \\
 & \texttt{rare\_words} & 405 & 78.5 & 80.7 & +2.2 \\
 & \texttt{csat\_socio} & 298 & 40.6 & 46.3 & +5.7 \\
 & \texttt{csat\_law} & 217 & 41.0 & 51.6 & +10.6 \\
 & \texttt{history} & 188 & 75.5 & 76.1 & +0.5 \\
 & \texttt{general\_knowledge} & 176 & 51.1 & 59.1 & +8.0 \\
 & \texttt{loan\_words} & 169 & 76.3 & 83.4 & +7.1 \\
 & \texttt{standard\_nomenclature} & 153 & 82.4 & 88.2 & +5.9 \\
 & \texttt{csat\_geo} & 150 & 54.0 & 70.7 & +16.7 \\
\bottomrule
\end{tabular}
\end{table}

\subsection{Interpretation and limitations}

\paragraph{No forgetting; accuracy rose.}
On these benchmarks we see no sign that task training damaged general Korean understanding.
Training on the decision format (the first-position identifier readout) improved answer
selection across multiple-choice questions in general. The gain is larger in Korean (4.6 to
6.9 points) than in English (3.7 points on MMLU, none on HellaSwag); 82\% of the training rows
were Korean (41,621 of 50,938, Table~\ref{tab:mixture}). We report the languages separately
and do not average them.

\paragraph{Part of the suite is close to the training distribution.}
KoBEST BoolQ and SentiNeg resemble in form the English BoolQ training rows and the NSMC
sentiment data used in training. Items whose text keys overlap the training mixture were
removed, but transfer of the task type is still mixed into these two results. KMMLU, HAE-RAE
and CLIcK are knowledge and culture questions with no similar data in training.

\paragraph{Date calculation remains weak.}
HAE-RAE date understanding stays near 50\% and does not improve. In the author-written probes
(Section~\ref{sec:probes}) the only questions \expone{} missed were also sum calculations.
Following the specification, calculation belongs in a code path, not in the model.

\paragraph{Calibration differs by benchmark.}
Applying the product temperature lowers the ECE of \expone{} on most benchmarks but raises it
on KoBEST and HellaSwag. We did not calibrate per benchmark or per task.

\paragraph{The readout baseline.}
The zero-shot score is zero-shot \emph{in this decision format}. It can be lower than what the
base model achieves with the benchmarks' official prompts or with generated answers. The
differences above are therefore the improvement obtained when the model is used through the
decision API, not evidence that the model's knowledge increased. For the same reason, these
first-position scores are not comparable with official leaderboard numbers, which use each
benchmark's official prompt and few-shot protocol.

\paragraph{Contamination.}
We did not measure whether the public base model saw these benchmarks during pretraining. The
upstream test splits are development data for Bobcat, and the results describe the change
caused by our training more reliably than the absolute level.
